% GENERATED FILE. Edit canonical lessons, then run npm run generate:books.
% canonical-source-hash: fnv1a64:37a5627f7fdf4636
% canonical-lessons: FA-C518-setar, FA-C518-daf, FA-C518-tabl, FA-C518-navakhtan, FA-C518-kuk-kardan

\chapter{Playing an Instrument}
\label{ch:fa-playing-an-instrument}
\hlchaptermodality{518}

\begin{chapteropening}
\textbf{By the end of this chapter:} \emph{I can talk about playing an instrument.}

Complete the last lesson of chapter 518: I can talk about playing an instrument.
\end{chapteropening}

\section[\texorpdfstring{\fa{سه} \fa{تار} (setâr)}{سه تار (setâr)}]{\fa{سه} \fa{تار} (setâr) — a setar (a three-stringed lute)}

\label{lesson:FA-C518-setar}



\begin{quote}
Before the new one: say the Persian for a santur, then the Persian for a tombak.
\end{quote}

\subsection*{You'll want to know: \fa{سه} \fa{تار}}
\textbf{\fa{سه} \fa{تار}} (\emph{setâr}) — \textquotedblleft{}a setar (a three-stringed lute)\textquotedblright{}.

\begin{quote}
\fa{او} \fa{شبها} \fa{برای} \fa{خودش} \fa{سه} \fa{تار} \fa{میزند}.
\end{quote}

\begin{quote}
\emph{u shabhâ barâ-ye khodash setâr mizanad.}
\end{quote}

\begin{quote}
At night she plays the setar for herself.
\end{quote}

\emph{se} 'three' + \emph{târ} 'string'. The Indian sitar is usually traced to it.

\begin{cousinweb}
\emph{setâr} is built on two native Persian words, \emph{se}, 'three', and \emph{târ}, 'string'. Hindi and Urdu \emph{sitār} are usually traced to it.
\end{cousinweb}

\begin{grammarlens}[title={What you've built}]
One more word for telling, explaining and describing at B1.
\end{grammarlens}

\subsection*{Your turn}
\begin{itemize}
\raggedright
  \item \emph{Say it:} \emph{setâr}
  \item \emph{Say it:} \emph{setâr}, once more
  \item \emph{Say it:} a sentence of your own with \emph{setâr}
  \item \emph{Recall:} say the Persian for a santur, then the Persian for a tombak, then say \emph{setâr} again
\end{itemize}

\begin{tcolorbox}[breakable,colback=teal!4,colframe=teal!35!black,title={Before you move on}]
What does \fa{سه} \fa{تار} mean? (\textquotedblleft{}A setar (a three-stringed lute)\textquotedblright{}.) Say it once more.
\end{tcolorbox}
\section[\texorpdfstring{\fa{دف} (daf)}{دف (daf)}]{\fa{دف} (daf) — a daf (a frame drum)}

\label{lesson:FA-C518-daf}



\begin{quote}
Before the new one: say the Persian for a tombak, then the Persian for a setar.
\end{quote}

\subsection*{You'll want to know: \fa{دف}}
\textbf{\fa{دف}} (\emph{daf}) — \textquotedblleft{}a daf (a frame drum)\textquotedblright{}.

\begin{quote}
\fa{زنها} \fa{در} \fa{جشن} \fa{دف} \fa{میزدند}.
\end{quote}

\begin{quote}
\emph{zanhâ dar jashn daf mizadand.}
\end{quote}

\begin{quote}
The women were playing frame drums at the celebration.
\end{quote}

\emph{daf zadan} — to play the daf.

\begin{cousinweb}
\emph{daf} is inherited: EtymDB traces it to Middle Persian \emph{daf} (EtymDB).
\end{cousinweb}

\begin{grammarlens}[title={What you've built}]
One more word for telling, explaining and describing at B1.
\end{grammarlens}

\subsection*{Your turn}
\begin{itemize}
\raggedright
  \item \emph{Say it:} \emph{daf}
  \item \emph{Say it:} \emph{daf}, once more
  \item \emph{Say it:} a sentence of your own with \emph{daf}
  \item \emph{Recall:} say the Persian for a tombak, then the Persian for a setar, then say \emph{daf} again
\end{itemize}

\begin{tcolorbox}[breakable,colback=teal!4,colframe=teal!35!black,title={Before you move on}]
What does \fa{دف} mean? (\textquotedblleft{}A daf (a frame drum)\textquotedblright{}.) Say it once more.
\end{tcolorbox}
\section[\texorpdfstring{\fa{طبل} (tabl)}{طبل (tabl)}]{\fa{طبل} (tabl) — a drum}

\label{lesson:FA-C518-tabl}



\begin{quote}
Before the new one: say the Persian for a setar, then the Persian for a daf.
\end{quote}

\subsection*{You'll want to know: \fa{طبل}}
\textbf{\fa{طبل}} (\emph{tabl}) — \textquotedblleft{}a drum\textquotedblright{}.

\begin{quote}
\fa{پسر} \fa{کوچک} \fa{طبل} \fa{میزند} \fa{و} \fa{همه} \fa{را} \fa{بیدار} \fa{میکند}.
\end{quote}

\begin{quote}
\emph{pesar-e kuchek tabl mizanad va hame râ bidâr mikonad.}
\end{quote}

\begin{quote}
The little boy bangs a drum and wakes everyone up.
\end{quote}

\emph{tabl zadan} — to beat a drum.

\begin{cousinweb}
\emph{tabl} comes from Arabic \emph{ṭabl}, 'drum' (EtymDB); Hindi \emph{tabla} goes back to the same word.
\end{cousinweb}

\begin{grammarlens}[title={What you've built}]
One more word for telling, explaining and describing at B1.
\end{grammarlens}

\subsection*{Your turn}
\begin{itemize}
\raggedright
  \item \emph{Say it:} \emph{tabl}
  \item \emph{Say it:} \emph{tabl}, once more
  \item \emph{Say it:} a sentence of your own with \emph{tabl}
  \item \emph{Recall:} say the Persian for a setar, then the Persian for a daf, then say \emph{tabl} again
\end{itemize}

\begin{tcolorbox}[breakable,colback=teal!4,colframe=teal!35!black,title={Before you move on}]
What does \fa{طبل} mean? (\textquotedblleft{}A drum\textquotedblright{}.) Say it once more.
\end{tcolorbox}
\section[\texorpdfstring{\fa{نواختن} (navâkhtan)}{نواختن (navâkhtan)}]{\fa{نواختن} (navâkhtan) — to play (an instrument)}

\label{lesson:FA-C518-navakhtan}



\begin{quote}
Before the new one: say the Persian for a daf, then the Persian for a drum.
\end{quote}

\subsection*{You'll want to know: \fa{نواختن}}
\textbf{\fa{نواختن}} (\emph{navâkhtan}) — \textquotedblleft{}to play (an instrument)\textquotedblright{}.

\begin{quote}
\fa{او} \fa{پیانو} \fa{را} \fa{خیلی} \fa{زیبا} \fa{مینوازد}.
\end{quote}

\begin{quote}
\emph{u piyâno râ kheyli zibâ minavâzad.}
\end{quote}

\begin{quote}
She plays the piano beautifully.
\end{quote}

Formal; present stem \emph{navâz-}. In speech people say \emph{zadan}: \emph{piyâno mizanad}.

\begin{cousinweb}
\emph{navâkhtan} is a native Persian verb, with the present stem \emph{navâz-} seen in \emph{navâzesh}. EtymDB lists it without a source, so its older history is not settled.
\end{cousinweb}

\begin{grammarlens}[title={What you've built}]
One more word for telling, explaining and describing at B1.
\end{grammarlens}

\subsection*{Your turn}
\begin{itemize}
\raggedright
  \item \emph{Say it:} \emph{navâkhtan}
  \item \emph{Say it:} \emph{navâkhtan}, once more
  \item \emph{Say it:} a sentence of your own with \emph{navâkhtan}
  \item \emph{Recall:} say the Persian for a daf, then the Persian for a drum, then say \emph{navâkhtan} again
\end{itemize}

\begin{tcolorbox}[breakable,colback=teal!4,colframe=teal!35!black,title={Before you move on}]
What does \fa{نواختن} mean? (\textquotedblleft{}To play (an instrument)\textquotedblright{}.) Say it once more.
\end{tcolorbox}
\section[\texorpdfstring{\fa{کوک} \fa{کردن} (kuk kardan)}{کوک کردن (kuk kardan)}]{\fa{کوک} \fa{کردن} (kuk kardan) — to tune (an instrument); to wind (a clock)}

\label{lesson:FA-C518-kuk-kardan}



\begin{quote}
Before the new one: say the Persian for a drum, then the Persian for to play.
\end{quote}

\subsection*{You'll want to know: \fa{کوک} \fa{کردن}}
\textbf{\fa{کوک} \fa{کردن}} (\emph{kuk kardan}) — \textquotedblleft{}to tune (an instrument); to wind (a clock)\textquotedblright{}.

\begin{quote}
\fa{او} \fa{قبل} \fa{از} \fa{کنسرت} \fa{تارش} \fa{را} \fa{کوک} \fa{کرد}.
\end{quote}

\begin{quote}
\emph{u qabl az konsert târash râ kuk kard.}
\end{quote}

\begin{quote}
He tuned his tar before the concert.
\end{quote}

\emph{kuk} 'tuning' + \emph{kardan}. \emph{kuk budan} — to be in tune, and in speech 'to be in good form'.

\begin{cousinweb}
\emph{kuk} means 'tuning' and 'winding', but no source used here gives its history, and its origin is not settled. \emph{kardan} is inherited Persian.
\end{cousinweb}

\begin{grammarlens}[title={What you've built}]
One more word for telling, explaining and describing at B1.
\end{grammarlens}

\subsection*{Your turn}
\begin{itemize}
\raggedright
  \item \emph{Say it:} \emph{kuk kardan}
  \item \emph{Say it:} \emph{kuk kardan}, once more
  \item \emph{Say it:} a sentence of your own with \emph{kuk kardan}
  \item \emph{Recall:} say the Persian for a drum, then the Persian for to play, then say \emph{kuk kardan} again
\end{itemize}

\begin{tcolorbox}[breakable,colback=teal!4,colframe=teal!35!black,title={Before you move on}]
What does \fa{کوک} \fa{کردن} mean? (\textquotedblleft{}To tune (an instrument); to wind (a clock)\textquotedblright{}.) Say it once more.
\end{tcolorbox}
